\documentclass[11pt,a4paper]{article}

\usepackage{ctex}
\usepackage[T1]{fontenc}
\usepackage[top=2.5cm, bottom=2.5cm, left=2.5cm, right=2.5cm]{geometry}
\usepackage{booktabs}
\usepackage{array}
\usepackage{enumitem}
\usepackage{tabularx}
\usepackage{amssymb}
\usepackage{xcolor}
\definecolor{primary}{HTML}{1a1a2e}
\definecolor{accent}{HTML}{3b82f6}
\definecolor{success}{HTML}{22c55e}
\definecolor{danger}{HTML}{ef4444}
\definecolor{warning}{HTML}{f59e0b}
\definecolor{graytext}{HTML}{6b7280}
\definecolor{progress}{HTML}{8b5cf6}
\usepackage[hidelinks]{hyperref}
\hypersetup{colorlinks=true, linkcolor=primary, urlcolor=accent}
\usepackage{titlesec}
\titleformat{\section}{\Large\bfseries\color{primary}}{}{0em}{}[\titlerule]
\titleformat{\subsection}{\large\bfseries\color{primary}}{}{0em}{}
\usepackage{fancyhdr}
\pagestyle{fancy}
\fancyhf{}
\fancyhead[L]{\small\color{graytext} 企业级 AI Agent 应用商店 -- 需求文档}
\fancyhead[R]{\small\color{graytext} DEMAND-008}
\fancyfoot[C]{\small\color{graytext} \thepage}
\renewcommand{\headrulewidth}{0.4pt}

\title{\textbf{DEMAND-008：目录框架重构与实习清单融合}}
\author{万能产品小助手（PM AI）\\魏子政 审阅}
\date{2026-05-06 / 版本 v0.1}

\begin{document}

\maketitle

\begin{center}
\begin{tabular}{>{\bfseries}p{3cm} p{10cm}}
\toprule
提出日期 & 2026-05-06 \\
提出者 & 魏子政 \\
优先级 & \textcolor{danger}{P0} \\
状态 & \textcolor{progress}{进行中} \\
总需求编号 & D-014 \\
前置需求 & D-011（MVP）、D-012（前端重写）、D-013（流转修复） \\
参考文档 & \texttt{docs/intern-dev-checklist.tex}（实习生开发清单） \\
\bottomrule
\end{tabular}
\end{center}

\vspace{1em}

\section{需求背景}

魏子政要求将 MVP 阶段生成的代码清空，重新编排前后端目录框架。同时要求将 \texttt{docs/intern-dev-checklist.tex}（实习生开发清单）中的工程规范与现有 MASTER 总需求做融合，形成正式的、面向多人协作的目录结构。

\textbf{核心目标}：
\begin{enumerate}[nosep]
  \item 清空 MVP 代码，保留配置文件和目录骨架
  \item 基于实习清单的完整工程视角，重新设计前后端目录结构
  \item 将目录结构与 MASTER 需求做映射，明确每个目录/文件对应哪个需求
  \item 为后续 Cursor 代码生成提供清晰的目录蓝图
\end{enumerate}

\section{实习清单核心内容提取}

\texttt{intern-dev-checklist.tex} 是一份面向实习生的 14 章开发指南，覆盖以下领域：

\begin{table}[h]
\centering
\begin{tabularx}{\linewidth}{clX}
\toprule
\textbf{章} & \textbf{主题} & \textbf{与 MASTER 需求的映射} \\
\midrule
1 & 环境准备 & D-008（开发环境搭建） \\
2 & 前端开发（结构+任务+规范） & D-012（前端重写） \\
3 & 后端开发（结构+任务+规范） & D-011（MVP 后端） \\
4 & 数据库 & D-011（PostgreSQL + Alembic） \\
5 & 搜索引擎 & D-015（搜索，待建） \\
6 & 安全扫描 & D-011（三层扫描已完成） \\
7 & 异步任务 & D-011（BackgroundTasks） \\
8 & 对象存储 & D-011（MinIO） \\
9 & 认证授权 & D-011（JWT） \\
10 & 容器化部署 & D-016（生产部署，待建） \\
11 & 日志监控 & D-017（可选项） \\
12 & 开发路线图 & 整体迭代规划 \\
13 & 企业级演进方向 & 长期架构升级计划 \\
14 & 名词速查表 & 团队知识库 \\
\bottomrule
\end{tabularx}
\end{table}

\section{前端目录框架}

\subsection{设计原则}

\begin{itemize}[nosep]
  \item \textbf{admin / portal 分离}：管理后台与开发者门户分目录，职责清晰
  \item \textbf{composables}：可复用逻辑抽离（如 \texttt{useAuth}、\texttt{usePagination}）
  \item \textbf{API 层统一}：所有接口调用集中在 \texttt{src/api/}，组件不直接发请求
  \item \textbf{TypeScript 优先}：所有 \texttt{.vue} 和 \texttt{.ts} 文件，不用 \texttt{.js}
\end{itemize}

\subsection{目录结构}

\begin{verbatim}
frontend/
|-- src/
|   |-- views/                  <- 页面文件
|   |   |-- admin/              <- 管理后台
|   |   |   |-- apps/           <- 应用管理（上架/下架/删除）
|   |   |   |-- reviews/        <- 审批管理（审批工作台）
|   |   |   |-- users/          <- 用户管理（RBAC 角色分配）
|   |   |-- portal/             <- 开发者门户
|   |   |   |-- explore/        <- 应用市场/发现页
|   |   |   |-- submit/         <- 提交应用页
|   |   |   |-- profile/        <- 开发者主页（我的应用）
|   |   |-- LoginView.vue       <- 登录页
|   |   |-- RegisterView.vue    <- 注册页
|   |-- components/             <- 可复用组件
|   |   |-- ui/                 <- 基础 UI 组件（Element Plus 二次封装）
|   |   |-- AppCard.vue         <- 应用卡片（市场列表用）
|   |   |-- SearchBar.vue       <- 搜索栏（带防抖）
|   |   |-- ScanResult.vue      <- 扫描结果展示
|   |   |-- StatusTag.vue       <- 状态标签（scanning/pending/published/rejected）
|   |-- composables/            <- Vue 3 Composables（可复用逻辑）
|   |   |-- useAuth.ts          <- 认证状态
|   |   |-- usePagination.ts    <- 分页逻辑
|   |   |-- useSearch.ts        <- 搜索防抖
|   |-- stores/                 <- Pinia 状态管理
|   |   |-- auth.ts             <- 用户认证 store
|   |   |-- skills.ts           <- Skill 数据 store
|   |   |-- ui.ts               <- UI 状态（侧边栏、暗色模式等）
|   |-- api/                    <- API 调用封装（axios）
|   |   |-- client.ts           <- axios 实例 + 拦截器
|   |   |-- auth.ts             <- 登录/注册/Token 刷新
|   |   |-- skills.ts           <- Skill CRUD + 搜索 + 下载 + 安装
|   |   |-- reviews.ts          <- 审批接口
|   |   |-- users.ts            <- 用户管理接口
|   |   |-- workflow.ts         <- 工作流/状态流转
|   |   |-- types.ts            <- 全局 TypeScript 类型定义
|   |-- router/                 <- Vue Router 路由配置
|   |   |-- index.ts            <- 路由定义 + 路由守卫
|   |-- utils/                  <- 工具函数
|   |   |-- jwt.ts              <- JWT 解析
|   |   |-- format.ts           <- 日期/文件大小格式化
|   |-- styles/                 <- 全局样式
|   |   |-- variables.css       <- CSS 变量（主题色）
|   |   |-- global.css          <- 全局样式重置
|   |   |-- element-override.css <- Element Plus 主题覆盖
|   |-- App.vue                 <- 根组件
|   |-- main.ts                 <- 入口文件
|-- public/
|   |-- icons/                  <- 静态图标
|-- index.html
|-- vite.config.ts
|-- tsconfig.json
|-- package.json
\end{verbatim}

\subsection{目录与需求映射}

\begin{table}[h]
\centering
\begin{tabularx}{\linewidth}{llX}
\toprule
\textbf{目录/文件} & \textbf{对应需求} & \textbf{说明} \\
\midrule
views/portal/explore/ & D-015（搜索） & 应用市场 + 搜索 + 筛选 \\
views/portal/submit/ & D-011 & Skill 上传表单 \\
views/portal/profile/ & D-013 & 我的应用（所有状态） \\
views/admin/apps/ & D-011 & 管理员应用管理 \\
views/admin/reviews/ & D-011 & 审批工作台 \\
views/admin/users/ & D-018（RBAC） & 用户管理 + 角色分配 \\
components/ScanResult.vue & D-011 & 三层扫描结果展示 \\
composables/useSearch.ts & D-015 & 搜索防抖 + OpenSearch 对接 \\
api/skills.ts & D-013 & 下载 + 安装接口 \\
\bottomrule
\end{tabularx}
\end{table}

\section{后端目录框架}

\subsection{设计原则}

\begin{itemize}[nosep]
  \item \textbf{core/}：核心配置与安全模块，独立于业务
  \item \textbf{models/}：SQLAlchemy ORM 模型，每个实体一个文件
  \item \textbf{routers/}：API 路由按功能分文件，每个文件一个 \texttt{APIRouter}
  \item \textbf{services/}：业务逻辑层，Router 不直接操作数据库
  \item \textbf{schemas/}：Pydantic 数据模型，请求/响应格式定义
  \item \textbf{tasks/}：asyncio 后台任务
  \item \textbf{tests/}：单元测试 + 集成测试
\end{itemize}

\subsection{目录结构}

\begin{verbatim}
backend/
|-- app/
|   |-- main.py               <- FastAPI 入口，注册路由 + 启动事件
|   |-- config.py             <- 配置（环境变量，Pydantic Settings）
|   |-- database.py           <- 数据库连接 + Session 工厂
|   |-- core/                 <- 核心模块
|   |   |-- security.py       <- JWT 签发/验证 + 密码哈希
|   |   |-- deps.py           <- FastAPI 通用依赖（get_db, get_current_user）
|   |   |-- permissions.py    <- RBAC 权限装饰器
|   |-- models/               <- SQLAlchemy ORM 模型
|   |   |-- base.py           <- Base + 公共字段（id, created_at, updated_at）
|   |   |-- user.py           <- 用户表
|   |   |-- skill.py          <- Skill 表
|   |   |-- skill_version.py  <- Skill 版本表
|   |   |-- scan_result.py    <- 扫描结果表
|   |   |-- review.py         <- 审批记录表
|   |-- schemas/              <- Pydantic 数据模型
|   |   |-- auth.py           <- 登录/注册请求响应
|   |   |-- skill.py          <- Skill CRUD 请求响应
|   |   |-- review.py         <- 审批请求响应
|   |   |-- common.py         <- 分页、通用响应
|   |-- routers/              <- API 路由
|   |   |-- auth.py           <- /api/auth/*
|   |   |-- skills.py         <- /api/skills/*
|   |   |-- reviews.py        <- /api/reviews/*
|   |   |-- search.py         <- /api/search/*（OpenSearch 对接）
|   |   |-- users.py          <- /api/users/*（RBAC 管理）
|   |   |-- workflow.py       <- /api/workflow/*（状态流转）
|   |-- services/             <- 业务逻辑层
|   |   |-- auth_service.py   <- 认证业务逻辑
|   |   |-- skill_service.py  <- Skill CRUD 逻辑
|   |   |-- scan_service.py   <- 三层扫描编排
|   |   |-- semgrep_scan.py   <- Semgrep 扫描
|   |   |-- clamav_scan.py    <- ClamAV 扫描
|   |   |-- llm_scan.py       <- LLM 语义分析
|   |   |-- search_service.py <- OpenSearch 搜索逻辑
|   |   |-- skill_package_service.py <- 安装包管理（下载/安装到本地）
|   |   |-- seed.py           <- 种子数据
|   |-- tasks/                <- 后台任务
|   |   |-- scan_task.py      <- 扫描任务触发
|   |   |-- notify_task.py    <- 通知任务（预留）
|   |-- utils/                <- 工具函数
|   |   |-- minio_client.py   <- MinIO 对象存储客户端
|   |   |-- password.py       <- 密码哈希
|   |   |-- opensearch.py     <- OpenSearch 客户端
|-- tests/
|   |-- unit/                 <- 单元测试
|   |-- integration/          <- 集成测试
|-- alembic/                  <- 数据库迁移
|   |-- versions/
|-- requirements.txt
|-- Dockerfile
|-- .env                      <- 环境变量（不提交 Git）
\end{verbatim}

\subsection{目录与需求映射}

\begin{table}[h]
\centering
\begin{tabularx}{\linewidth}{llX}
\toprule
\textbf{目录/文件} & \textbf{对应需求} & \textbf{说明} \\
\midrule
core/security.py & D-011 & JWT + 密码哈希 \\
core/permissions.py & D-018（RBAC） & 角色权限控制 \\
models/ & D-008 & 所有数据库实体 \\
routers/search.py & D-015 & OpenSearch 搜索接口 \\
services/search\_service.py & D-015 & 搜索业务逻辑 \\
services/scan\_* & D-011 & 三层扫描（已完成） \\
services/skill\_package\_service.py & D-013 & 下载 + 安装到 OpenClaw \\
tasks/scan\_task.py & D-011 & 后台扫描任务 \\
utils/opensearch.py & D-015 & OpenSearch 客户端封装 \\
\bottomrule
\end{tabularx}
\end{table}

\section{编码规范（融合实习清单）}

\subsection{前端编码规范}

\begin{enumerate}[nosep]
  \item 组件用 TypeScript，文件名 PascalCase（如 \texttt{AppCard.vue}）
  \item 遵循 \texttt{<script setup lang="ts">} 风格
  \item UI 组件优先用 Element Plus，复杂交互可二次封装
  \item API 调用统一放 \texttt{src/api/}，用 axios 封装，不在组件里直接 \texttt{fetch}
  \item 状态管理用 Pinia，按功能模块分 store
  \item 可复用逻辑抽到 \texttt{composables/}
  \item 样式用 CSS 变量 + Element Plus 主题覆盖
  \item 颜色方案：黑白灰极简（主色 \texttt{\#1a1a2e}，辅助 \texttt{\#f5f5f5}，文字 \texttt{\#333333}）
\end{enumerate}

\subsection{后端编码规范}

\begin{enumerate}[nosep]
  \item 所有接口用 \texttt{async/await}
  \item 数据校验用 Pydantic，不手动校验
  \item 数据库操作放 \texttt{services/}，不在 Router 里直接操作
  \item 敏感信息放环境变量，不硬编码
  \item 每个 Router 加 \texttt{tags}，方便 Swagger 文档
  \item 每个模块有对应的 \texttt{tests/unit/} 和 \texttt{tests/integration/} 测试
\end{enumerate}

\section{迭代规划（融合实习清单路线图）}

\begin{table}[h]
\centering
\begin{tabularx}{\linewidth}{clX}
\toprule
\textbf{阶段} & \textbf{重点} & \textbf{对应需求} \\
\midrule
\textbf{V1.1} & 目录框架重构 + 环境重建 & D-014（本文档） \\
& 清空旧代码，按新目录结构重建 & \\
& 安装依赖，确认服务可启动 & \\
\midrule
\textbf{V1.2} & 核心功能重建 & D-011 + D-012 + D-013 \\
& Skill CRUD + 文件上传 + MinIO & \\
& 三层扫描 + 审批工作台 & \\
& 登录/注册 + JWT & \\
& 应用市场 + 详情页 + 我的应用 & \\
& 下载 zip + 安装到 OpenClaw & \\
\midrule
\textbf{V1.3} & 搜索 + RBAC & D-015 + D-018 \\
& OpenSearch 全文 + 语义搜索 & \\
& RBAC 四角色权限体系 & \\
& 用户管理页面 & \\
\midrule
\textbf{V1.4} & 打磨 + 部署 & D-016 \\
& UI 优化 + 响应式 & \\
& Docker 生产配置 & \\
& 健康检查 + 日志 & \\
\bottomrule
\end{tabularx}
\end{table}

\section{企业级演进方向（引用实习清单第 13 章）}

以下为 MVP 后的升级计划，不做当前迭代，但目录结构需预留扩展空间：

\begin{table}[h]
\centering
\begin{tabularx}{\linewidth}{llX}
\toprule
\textbf{维度} & \textbf{现在（MVP）} & \textbf{未来（企业级）} \\
\midrule
认证 & 自建 JWT & Keycloak（LDAP/SSO） \\
工作流 & 数据库状态机 & Temporal（分布式工作流） \\
消息 & asyncio & Kafka（消息队列） \\
存储 & PostgreSQL JSONB & PostgreSQL + MongoDB \\
部署 & Docker Compose & Kubernetes + Helm \\
监控 & structlog & Prometheus + Grafana + Loki \\
前端 & Element Plus & 自研组件库 \\
\bottomrule
\end{tabularx}
\end{table}

\section{验收标准}

\begin{itemize}
  \item[$\square$] 前端目录结构符合上述规范
  \item[$\square$] 后端目录结构符合上述规范
  \item[$\square$] MVP 旧代码已清空（仅保留 .gitkeep 和 \texttt{\_\_init\_\_.py}）
  \item[$\square$] \texttt{npm install} + \texttt{npm run build} 无报错（即使页面为空）
  \item[$\square$] 后端 \texttt{pip install} + \texttt{uvicorn} 可启动（\texttt{/health} 正常）
  \item[$\square$] 需求文档 MASTER 已同步更新
  \item[$\square$] DOC-NAMING-CONVENTION 索引已更新
\end{itemize}

\section*{更新记录}

\begin{tabular}{lll}
\toprule
版本 & 日期 & 变更内容 \\
\midrule
v0.1 & 2026-05-06 & 初始版本 \\
\bottomrule
\end{tabular}

\end{document}
